\section{DP $O(n\sqrt n)$ por cota triangular}
\whenuse{Un par\'ametro del estado crece o decrece estrictamente en cada salto,
por lo que la suma m\'inima de los saltos limita su cantidad a $O(\sqrt n)$.}

Usar cuando un estado guarda una longitud/salto que cambia estrictamente en
cada transici\'on. No pueden existir demasiados cambios: $1+2+\cdots+t\le n$
implica $t=O(\sqrt n)$. Guardar s\'olo estados hasta ese l\'imite suele bajar una
DP aparentemente cuadr\'atica a $O(n\sqrt n)$.

\complexity{tiempo $O(n\sqrt n)$; memoria $O(n\sqrt n)$ o $O(\sqrt n)$ con
rolling arrays.}

\begin{lstlisting}[style=cpp]
int triangular_state_limit(int n) {
    int length = 0;
    while (1LL * (length + 1) * (length + 2) / 2 <= n) ++length;
    return length;
}

// Esqueleto: adaptar la transicion al problema.
vector<vector<ll>> triangular_bound_dp(int n) {
    int limit = triangular_state_limit(n);
    const ll infinity = (1LL << 60);
    vector<vector<ll>> dp(n + 1, vector<ll>(limit + 1, infinity));
    dp[0][0] = 0;
    for (int position = 0; position < n; ++position)
        for (int length = 0; length <= limit; ++length) if (dp[position][length] < infinity) {
            dp[position + 1][length] = min(dp[position + 1][length], dp[position][length]);
            int jump = length + 1;
            if (position + jump <= n && jump <= limit)
                dp[position + jump][jump] = min(dp[position + jump][jump],
                                                dp[position][length] + 1);
        }
    return dp;
}
\end{lstlisting}

\warning{La cota debe demostrarse con la suma m\'inima de los saltos; no basta
con observar que aparece $\sqrt n$ en los l\'imites.}
